\documentclass[11pt,a4paper]{article}
\usepackage[T2A]{fontenc}
\usepackage[utf8]{inputenc}
\usepackage[russian,english]{babel}
\usepackage{amsmath,amssymb,amsthm,mathtools}
\usepackage{setspace}
\usepackage{enumitem}
\usepackage[hidelinks]{hyperref}
\setstretch{1.3}
\frenchspacing
\hyphenpenalty=10000

\newcommand{\head}[1]{\par\noindent\textbf{\underline{#1}}\quad}
\newcommand{\sheettitle}[4]{% title, subtitle, homework-flag, date
  \begin{center}
  {\huge\bfseries #1\par}\vspace{\dimexpr1.2em-1pt\relax minus 3.6pt}
  {\Large #2\par}\vspace{\dimexpr1em+1pt\relax minus 1.2pt}
  \if\relax\detokenize{#3}\relax\else{\Large #3\par}\vspace{\dimexpr1.15em-1pt\relax minus 3.6pt}\fi
  {\foreignlanguage{russian}{#4}\par}
  \end{center}\vspace{0.5pt minus 1.2pt}}
\newcommand{\src}[1]{\unskip\nobreak\hfil\penalty50\hskip1em\hbox{}\nobreak\hfill(#1){\parfillskip=0pt\par}}
\DeclareMathOperator{\tr}{tr}
\DeclareMathOperator{\adj}{adj}
\newcommand{\T}{\mathsf T}
\begin{document}
\sheettitle{Linear algebra -- 1}{Determinants and invertibility}{Homework}{12 сентября 2026 года}

\head{Theoretical minimum.} The adjugate satisfies $A\adj A=(\det A)E$ and has rank $n$, $1$ or
$0$ according to the rank of $A$; Jacobi's formula differentiates a determinant and gives
$\det e^A=e^{\tr A}$. A determinant of a matrix which is a small-rank perturbation of a simple one
is computed by $\det(A+UV^\T)=\det A\cdot\det(E+V^\T A^{-1}U)$. An integer matrix is invertible
over $\mathbb Z$ exactly when $\det A=\pm1$, and since $\det$ is a polynomial with integer
coefficients, identities transfer to any commutative ring and reduction modulo a prime proves
non-vanishing. Singularity has three equivalent faces --- dependent columns, a zero eigenvalue, a
non-trivial kernel --- and the art is to pick the right one. Finally, over $\mathbb R$ a
determinant is a continuous function of the entries, so a determinant identity proved on a dense
set holds everywhere: this is how one passes from invertible matrices to all of them.

\par\vspace{-1pt}\noindent\rule{\linewidth}{0.4pt}\par\vspace{7pt}

\begin{enumerate}
\item (VJIMC 2013) Let $A=(a_{ij})$ and $B=(b_{ij})$ be two real $10\times10$ matrices such that
$a_{ij}=b_{ij}+1$ for all $i,j$ and $A^3=O$. Prove that $\det B=0$.

\item (IMC 2015) For any integer $n\geqslant2$ and two $n\times n$ matrices with real entries
$A,B$ satisfying $A^{-1}+B^{-1}=(A+B)^{-1}$, prove that $\det A=\det B$. Does the same
conclusion hold for matrices with complex entries?

\item (IMC 2014) Determine all pairs $(a,b)$ of real numbers for which there exists a unique
symmetric $2\times2$ matrix $M$ with real entries satisfying $\tr M=a$ and $\det M=b$.

\item (Romania national olympiad 2017) Let $n\geqslant2$ and $k$ be natural numbers and let
$A_1,\dots,A_k$ be symmetric real $n\times n$ matrices. Prove that the following are equivalent:
\begin{enumerate}[label=\arabic*),align=left]
\item $\det\bigl(\sum_{i=1}^kA_i^2\bigr)=0$;
\item $\det\bigl(\sum_{i=1}^kA_iB_i\bigr)=0$ for all $B_1,\dots,B_k\in M_n(\mathbb R)$.
\end{enumerate}

\item (Romania national olympiad 1995) Let $n,k$ be positive integers and let
$A_i\in M_n(\mathbb R)$, $i=1,\dots,k$. Prove that
$\det\bigl(\sum_{i=1}^kA_i^\T A_i\bigr)\geqslant0$.

\item (Romania national olympiad 2011) Let $A$ be a singular real matrix of order $n\geqslant2$
and let $\adj A$ be its adjugate. Prove that $\tr(\adj A)\ne-1$ if and only if the matrix
$E+\adj A$ is invertible.

\item (Romania national olympiad 2008) Let $A$ be an $n\times n$ matrix with complex entries.
Prove that $A^{-1}=\overline A$ if and only if there exists an invertible complex matrix $B$ with
$A=B^{-1}\overline B$.

\item (Romania national olympiad 2025) Let $n$ be a positive integer and let $a,b$ be complex
numbers with $a\ne0$ and $b^k\ne1$ for every $k\in\{1,2,\dots,n\}$. The matrices
$A,B\in M_n(\mathbb C)$ satisfy $BA=aE+bAB$. Prove that $A$ and $B$ are invertible.

\item (Romania national olympiad 2012) Let $n\geqslant2$ and $1\leqslant k\leqslant n-1$. Prove
that if $A\in M_n(\mathbb C)$ has exactly $k$ minors of order $n-1$ equal to zero, then
$\det A\ne0$.

\item (Romania national olympiad 2014) Let $A\in M_4(\mathbb R)$ be invertible with
$\tr A=\tr(\adj A)\ne0$. Show that $A^2+E$ is singular if and only if there is a non-zero matrix
in $M_4(\mathbb R)$ anticommuting with $A$.

\item (VJIMC 1991) Compute the determinant of an antisymmetric matrix $(a_{ij})_{i,j=1}^n$,
where $a_{ij}=-a_{ji}$ and $a_{ii}=0$, for odd $n$, and explain what changes for even $n$.

\item (\foreignlanguage{russian}{Мехмат}, 2007) Prove Carleman's inequality: for every
real matrix~$A=(a_{ij})_{i,j\leqslant n}$
\[
|\det(E+A)|\leqslant\exp\Bigl(\tr A+\tfrac12\sum_{i,j=1}^na_{ij}^2\Bigr).
\]

\item Prove that if all entries of an integer $n\times n$ matrix are odd, then $2^{n-1}$ divides
its determinant, and show that the exponent $n-1$ cannot be improved.

\item (Putnam 2008) Alan and Barbara play a game in which they take turns filling entries of an
initially empty $2008\times2008$ array. Alan plays first. At each turn a player chooses a real
number and places it in a vacant entry. The game ends when all the entries are filled. Alan wins
if the determinant of the resulting matrix is non-zero, Barbara wins if it is zero. Which player
has a winning strategy?

\item (Putnam 2002) In Determinant Tic-Tac-Toe, Player 1 enters a 1 in an empty $3\times3$
matrix, Player 0 counters with a 0 in a vacant position, and play continues in turn until the
matrix is filled with five ones and four zeros. Player 0 wins if the determinant is 0 and Player 1
wins otherwise. Assuming optimal play, who wins and how?

\item (CIIM 2009) Determine whether for every natural $n$ there is a real $n\times n$ matrix whose
determinant is zero but such that changing any single entry produces a matrix with non-zero
determinant.
\end{enumerate}
\end{document}
